\begin{explanationbox}
\textbf{一、先识别问题：从“角”回到“边”。}

焦点 $F_1,F_2$ 在同一条 $x$ 轴上，角平分线从 $P$ 射向底边
$F_1F_2$。因此最自然的第一反应不是设斜率、不是求三角函数，
而是\textbf{角平分线定理}：
\[
\frac{F_1M}{MF_2}=\frac{PF_1}{PF_2}.
\]
设 $PF_1=r_1,\ PF_2=r_2$，由于 $F_1M=c+m$、$MF_2=c-m$，
\[
\frac{c+m}{c-m}=\frac{r_1}{r_2}
\quad\Longrightarrow\quad
\boxed{m=c\frac{r_1-r_2}{r_1+r_2}}.\tag{A}
\]
至此，\textbf{已经把“求角平分线交点”变成“处理两条焦距”}。

\medskip
\textbf{二、为什么自然想到“平方差”？}

因为两条焦距的坐标公式只有焦点横坐标 $\pm c$ 不同：
\[
\begin{aligned}
r_1^2&=(x_0+c)^2+y_0^2,\\
r_2^2&=(x_0-c)^2+y_0^2.
\end{aligned}
\]
相减能直接消去 $x_0^2,y_0^2,c^2$：
\[
\boxed{(r_1-r_2)(r_1+r_2)=4cx_0}.\tag{B}
\]
从 (A) 看见\textbf{差与和的商}，从 (B) 看见\textbf{差与和的积}，
两者配合即可只留下焦距差的平方。

\medskip
\textbf{三、为什么结果只有 $a^2$ 而没有 $b,c$？}

双曲线的本质是\textbf{到两个焦点的距离差的绝对值恒为 $2a$}：
\[
|PF_1-PF_2|=|r_1-r_2|=2a.
\]
所以 $b,c$ 在计算过程中消掉：
\[
mx_0=\frac{(r_1-r_2)^2}{4}=a^2.
\]
这比死记“截距公式”更可靠：它是\emph{角平分线定理、距离平方差、双曲线第一定义}的自然结果。

\medskip
\textbf{四、从“算出 $M$”提升到“认出切线”。}

双曲线的切线不是凭空出现的。焦距差固定，
沿曲线作很小的位移时，两条焦距的变化相互抵消。
在几何上，这使\textbf{焦点三角形顶角的内角平分线恰好成为切线方向}。
因此只要认出“$PF_1$、$PF_2$ 和顶角内分线”，就应该联想到
\[
PM\ \text{是切线},\qquad \frac{x_0x}{a^2}-\frac{y_0y}{b^2}=1.
\]
\textbf{对照提醒：}椭圆把“两焦距之和”固定，其焦点三角形顶角的
\emph{内角平分线是法线方向}；双曲线固定“两焦距之差”，内角平分线则是\emph{切线方向}。

\medskip
\textbf{五、内切圆帮助记住的另一个“固定点”。}

作 $\triangle PF_1F_2$ 的内切圆，切 $F_1F_2$ 于 $T$。
从同一点向圆引的切线长度相等，故
\[
F_1T-F_2T=PF_1-PF_2=\pm2a.
\]
又有 $F_1T+F_2T=2c$，得到切点 $T$ 在该支顶点。
\textbf{注意：} $T$ 是\emph{圆的切点}，$M$ 是\emph{角平分线与轴的交点}；
前者不随 $P$ 改变，后者随 $x_0$ 变化。不能把 $T$ 和 $M$ 混为一谈。

\medskip
\textbf{考试识别信号：}
\begin{itemize}
  \item 题目给出双曲线、两个焦点以及 $\angle F_1PF_2$ 的\textbf{内角平分线}；
  \item 所求为角平分线与焦点轴交点、横截距或相关直线方程；
  \item 或者问\textbf{焦点三角形的内切圆}与焦点轴相切于哪里。
\end{itemize}
先检查 $y_0\ne0$，再按“\textbf{角分线--焦距比--平方差--$2a$}”处理。
\end{explanationbox}
